% Figure: comparison of drag-free, linear-drag, and quadratic-drag
% trajectories from the same launch condition, computed numerically.
% Curves are precomputed sample points (no closed form for drag cases).
\begin{figure}[htbp]
\centering
\begin{tikzpicture}
\begin{axis}[
  width=12cm, height=7cm,
  xlabel={horizontal distance $x$ (\si{\meter})},
  ylabel={height $y$ (\si{\meter})},
  xmin=0, xmax=95, ymin=0, ymax=26,
  axis lines=left,
  legend pos=north east,
  legend cell align=left,
  grid=major, grid style={enggray!20},
  tick label style={font=\scriptsize},
  label style={font=\small},
  legend style={font=\scriptsize},
]
% Drag-free: v0=30, theta=45, coeff 0.0109 (from previous figure).
\addplot[engblue, very thick, domain=0:91.7, samples=120] {1.0*x - 0.0109*x^2};
\addlegendentry{no drag}
% Linear drag, sampled from projectile_drag.py (K_LIN = 0.02 N.s/m).
\addplot[engamber, very thick, dashed] coordinates {
  (0,0) (4.1,3.9) (7.9,7.2) (11.5,9.8) (14.8,11.9) (17.9,13.5)
  (20.8,14.6) (23.5,15.2) (26.0,15.5) (28.3,15.3) (30.5,14.8)
  (32.5,13.9) (34.4,12.7) (36.2,11.2) (37.8,9.4) (39.4,7.4)
  (40.8,5.1) (42.1,2.6) (43.3,0)};
\addlegendentry{linear drag}
% Quadratic drag, sampled from projectile_drag.py (tennis ball).
\addplot[engred, very thick] coordinates {
  (0,0) (4.0,3.8) (7.5,6.8) (10.7,9.1) (13.7,10.9) (16.4,12.2)
  (19.0,13.0) (21.4,13.4) (23.7,13.4) (25.9,13.0) (28.0,12.2)
  (30.0,11.1) (31.8,9.7) (33.6,8.0) (35.3,6.0) (36.8,3.7)
  (38.3,1.2) (38.9,0)};
\addlegendentry{quadratic drag}
\end{axis}
\end{tikzpicture}
\caption[Effect of drag on a projectile trajectory]{Trajectory of a
projectile launched at $45^{\circ}$ with $v_{0}=30\,\si{\meter\per
\second}$, computed under three force models by the solver in
\texttt{projectile\_drag.py}. The drag-free parabola (blue) is
symmetric about its apex and reaches $91.7\,\si{\meter}$. Quadratic
drag $F=-b\,|\vec{v}|\vec{v}$ (red), the correct model for a tennis
ball in this speed range ($\mathrm{Re}\sim 10^{4}$ to $10^{5}$), cuts
the range to $38.9\,\si{\meter}$ and breaks the symmetry: the
descending branch is steeper than the ascending branch because the
body sheds horizontal speed throughout the flight. Linear drag (amber
dashed), valid only in the creeping-flow regime, lands at
$43.3\,\si{\meter}$; it is the wrong model for a ball of this size and
speed, included to show how far the regime assumption can move the
answer. The asymmetry is the visible signature of drag in the
project's video data.}
\label{fig:vol03ch04:projectile-drag}
\end{figure}
